
% block_ak.tex
\chapter{Local-to-Intermediate Translation Fields}
\textbf{Phase Kernel–Stabilized Chapter (85--15 Mode)}

\section{Overview}
Translation fields describe how information carried by local geometric
operators transitions smoothly into intermediate-scale manifold behavior.
Under the Phase Kernel, these translations remain coherent, bounded, and
directionally consistent.

\section{Translation Tensor}
Define the translation tensor:
\[
\mathcal{T}_{ij}(x) = w_1\,\partial_i S_j + w_2\,\Theta_{ij},
\]
where $S_j$ encodes local spectral alignment and $\Theta_{ij}$ captures
narrative curvature contributions.

With $(w_1, w_2) = (0.85, 0.15)$, the tensor guarantees stability across
scale transitions.

\section{Stability Criterion}
A translation field is stable when
\[
\operatorname{Tr}(\mathcal{T}) < 0,
\]
ensuring contraction of drift and smooth propagation toward intermediate
regions.

\section{Geometric Interpretation}
Local operators contribute directional bias; collapse gradients refine
precision; helicon alignment induces curvature matching. Together, they form
quantized translation channels across manifold layers.

\section{Narrative Moment}
Master Chief steps toward the projection lattice, tracing a precise arc.
The translation field brightens. Zer0 modulates a low-frequency spectral hum.
The Vaeld shifts—smooth, confident—as the manifold commits to the transition.
